\section{Introduction}
A number of studies to date have conducted climate-chemistry simulations, to investigate the plume overlap effect \cite{}. Through modelling of local plume dispersion and the ensuing chemical and microphysical transformations, they have found a general trend towards reducing \ce{NOx}- and contrail-related warming under increasing emissions concentrations. This is due to saturation effects which occur in high-emission airspace, as elaborated on in \ref{saturation}. Such studies do not however, provide far-reaching conclusions on the range of saturation effects that occur, under different formation flight configurations and atmospheric conditions. In chapter \ref{chapter3}, the practicality of controlled plume intersection was investigated, through real-world air traffic data analysis. Under the "significant plume encounter (SPE)" criteria laid out, it was found that nearly 40\% of all aircraft plumes in the high-density hotspots over mainland US and Europe, exhibited some degree of plume overlap. Chapter \ref{chapter4} approaches the problem from a different perspective, focusing primarily on the chemical saturation effects that would result from high \ce{NOx} emissions. A metric called \ce{NO2} threshold was developed, to represent the propensity of a region for \ce{NOx} saturation in the upper troposphere. The chapter concluded with a first-order photochemical box model experiment, which went some way towards quantifying the chemical saturation effect in terms of formation flight and plume intersection. The motivation behind this research was to begin to identify suitable regions of the atmosphere to exploit chemical saturation effects resulting from formation flight and plume overlap. 

Chapters \ref{chapter5} and \ref{chapter6} serve as a direct continuation of this body of research, aiming to provide further real world context through higher-fidelity aircraft plume modelling and local atmospheric chemistry analysis. The literature survey conducted throughout the duration of this PhD revealed a lack of pre-existing modelling architectures, suitable for simulating the emissions, dispersion and local chemical and physical evolution of overlapping aircraft exhaust plumes. Therefore it was deemed imperative to construct a novel framework to fit these requirements. This led to the development of the Gridded Plume Analysis Tool (GPAT). The following sections will delineate the underlying models, input parameters and datasets contained within the GPAT architecture.

